\documentclass[12pt,a4paper]{article}

\usepackage[a4paper,margin=1in]{geometry}
\usepackage{amsmath,amssymb}
\usepackage{graphicx}
\usepackage{booktabs}
\usepackage{float}
\usepackage{listings}
\usepackage{xcolor}
\usepackage{titlesec}
\usepackage{setspace}

\definecolor{codegreen}{rgb}{0,0.6,0}
\definecolor{codegray}{rgb}{0.5,0.5,0.5}
\definecolor{codepurple}{rgb}{0.58,0,0.82}

\lstset{
    language=Python,
    basicstyle=\ttfamily\small,
    keywordstyle=\color{blue},
    commentstyle=\color{codegreen},
    stringstyle=\color{red},
    numbers=left,
    numberstyle=\tiny\color{codegray},
    breaklines=true,
    frame=single,
    showstringspaces=false
}

\title{\textbf{DCT-Based JPEG Compression Study}\\
Detailed Code Explanation and Analysis}
\author{}
\date{}

\begin{document}

\maketitle

\tableofcontents
\newpage

\section{Introduction}

This project implements a simplified JPEG-style image compression system using the Discrete Cosine Transform (DCT). The objective is to study how JPEG compression works internally by applying frequency-domain transformation, quantization, and reconstruction on grayscale images.

The project demonstrates:
\begin{itemize}
    \item Frequency-domain image processing
    \item Lossy image compression
    \item Quantization-based information reduction
    \item Image quality analysis using PSNR and SSIM
    \item Energy compaction property of DCT
\end{itemize}

\section{Library Imports}

\begin{lstlisting}
import cv2
import matplotlib.pyplot as plt
import numpy as np
from scipy.fftpack import dctn, idctn
from skimage.metrics import structural_similarity as ssim
\end{lstlisting}

\subsection{Explanation}

\begin{itemize}
    \item \textbf{OpenCV (cv2)} is used for image loading and processing.
    \item \textbf{NumPy} is used for matrix operations and numerical computation.
    \item \textbf{Matplotlib} is used for image and graph visualization.
    \item \textbf{SciPy FFTPack} provides efficient implementations of DCT and inverse DCT.
    \item \textbf{scikit-image} provides the Structural Similarity Index (SSIM) metric.
\end{itemize}

\section{JPEG Quantization Matrix}

\begin{lstlisting}
BASE_QUANTIZATION_MATRIX = np.array([
    [16, 11, 10, 16, 24, 40, 51, 61],
    [12, 12, 14, 19, 26, 58, 60, 55],
    [14, 13, 16, 24, 40, 57, 69, 56],
    [14, 17, 22, 29, 51, 87, 80, 62],
    [18, 22, 37, 56, 68, 109, 103, 77],
    [24, 35, 55, 64, 81, 104, 113, 92],
    [49, 64, 78, 87, 103, 121, 120, 101],
    [72, 92, 95, 98, 112, 100, 103, 99]
], dtype=np.float32)
\end{lstlisting}

\subsection{Purpose}

This matrix controls how aggressively each DCT frequency coefficient is compressed.

\subsection{Why This Matrix Is Structured This Way}

\begin{itemize}
    \item Small values in the top-left preserve low-frequency coefficients.
    \item Large values in the bottom-right heavily compress high-frequency coefficients.
\end{itemize}

This works because human vision is more sensitive to low-frequency information and less sensitive to fine high-frequency details.

\section{Quality-Based Quantization Scaling}

\begin{lstlisting}
def get_quantization_matrix(quality=50):
    quality = np.clip(quality, 1, 100)
    scale = (5000 / quality) if quality < 50 else (200 - 2 * quality)
    return np.clip((BASE_QUANTIZATION_MATRIX * scale + 50) / 100, 1, 255)
\end{lstlisting}

\subsection{Purpose}

This function dynamically scales the JPEG quantization matrix based on a user-selected quality factor.

\subsection{Quality Factor}

\begin{itemize}
    \item Low quality value $\rightarrow$ stronger compression
    \item High quality value $\rightarrow$ better image quality
\end{itemize}

\subsection{Why Two Different Scaling Equations?}

JPEG quality scaling is nonlinear.

\[
scale = \frac{5000}{Q}, \quad Q < 50
\]

\[
scale = 200 - 2Q, \quad Q \geq 50
\]

Lower quality values produce aggressive compression while higher quality values preserve more information.

\section{DCT Compression Function}

\begin{lstlisting}
def process_image_dct(image, quant_matrix):
    h, w = image.shape
    img_float = np.float32(image) - 128

    blocks = img_float.reshape(h // 8, 8, w // 8, 8).transpose(0, 2, 1, 3)

    dct_blocks = dctn(blocks, axes=(-2, -1), norm='ortho')
    quantized = np.round(dct_blocks / quant_matrix)
    dequantized = quantized * quant_matrix
    recon = idctn(dequantized, axes=(-2, -1), norm='ortho')

    result = recon.transpose(0, 2, 1, 3).reshape(h, w) + 128
    return np.clip(result, 0, 255).astype(np.uint8)
\end{lstlisting}

\subsection{Step 1: Image Dimensions}

\begin{lstlisting}
h, w = image.shape
\end{lstlisting}

Gets image height and width.

\subsection{Step 2: Level Shift}

\begin{lstlisting}
img_float = np.float32(image) - 128
\end{lstlisting}

Pixel values are shifted from:

\[
[0,255] \rightarrow [-128,127]
\]

This centers data around zero, improving DCT efficiency.

\subsection{Step 3: Block Splitting}

\begin{lstlisting}
blocks = img_float.reshape(h // 8, 8, w // 8, 8).transpose(0, 2, 1, 3)
\end{lstlisting}

JPEG processes images in \(8 \times 8\) blocks because:
\begin{itemize}
    \item computationally efficient,
    \item good balance between quality and compression,
    \item minimizes visual artifacts.
\end{itemize}

\subsection{Why transpose()?}

The transpose reorganizes axes into:

\[
(block\_row, block\_col, 8, 8)
\]

allowing simultaneous vectorized block processing.

\subsection{Step 4: Forward DCT}

\begin{lstlisting}
dct_blocks = dctn(blocks, axes=(-2, -1), norm='ortho')
\end{lstlisting}

Transforms each \(8 \times 8\) block into frequency coefficients.

\subsection{DCT Formula}

\[
F(u,v)=\frac{1}{4}C(u)C(v)\sum_{x=0}^{7}\sum_{y=0}^{7}
f(x,y)\cos\left[\frac{(2x+1)u\pi}{16}\right]
\cos\left[\frac{(2y+1)v\pi}{16}\right]
\]

\subsection{Why DCT?}

DCT concentrates most image energy into low-frequency coefficients.

\subsection{Step 5: Quantization}

\begin{lstlisting}
quantized = np.round(dct_blocks / quant_matrix)
\end{lstlisting}

This is the lossy compression step.

\subsection{Quantization Formula}

\[
Q(u,v)=\text{round}\left(\frac{F(u,v)}{T(u,v)}\right)
\]

High-frequency coefficients become zero after quantization.

\subsection{Step 6: Dequantization}

\begin{lstlisting}
dequantized = quantized * quant_matrix
\end{lstlisting}

Approximate reconstruction of DCT coefficients.

\subsection{Step 7: Inverse DCT}

\begin{lstlisting}
recon = idctn(dequantized, axes=(-2, -1), norm='ortho')
\end{lstlisting}

Transforms frequency coefficients back into pixel values.

\subsection{Step 8: Reconstruct Image}

\begin{lstlisting}
result = recon.transpose(0, 2, 1, 3).reshape(h, w) + 128
\end{lstlisting}

Rebuilds the original image layout and reverses the earlier level shift.

\subsection{Step 9: Clipping}

\begin{lstlisting}
np.clip(result, 0, 255).astype(np.uint8)
\end{lstlisting}

Ensures all pixel values remain within the valid image range.

\section{Quality Metrics}

\begin{lstlisting}
def compute_metrics(original, compressed):
    mse = np.mean((original.astype(np.float32)
           - compressed.astype(np.float32)) ** 2)

    psnr = 100.0 if mse == 0 else \
           20 * np.log10(255 / np.sqrt(mse))

    ssim_score = ssim(original, compressed)

    return mse, psnr, ssim_score
\end{lstlisting}

\subsection{Mean Squared Error (MSE)}

\[
MSE = \frac{1}{MN}\sum_{x=1}^{M}\sum_{y=1}^{N}
[I(x,y)-K(x,y)]^2
\]

Measures average pixel error.

Lower MSE indicates better reconstruction quality.

\subsection{Peak Signal-to-Noise Ratio (PSNR)}

\[
PSNR = 20\log_{10}\left(\frac{255}{\sqrt{MSE}}\right)
\]

Measures signal quality in decibels.

Higher PSNR indicates better quality.

\subsection{Structural Similarity Index (SSIM)}

SSIM measures:
\begin{itemize}
    \item luminance similarity,
    \item contrast similarity,
    \item structural similarity.
\end{itemize}

SSIM range:

\[
0 \leq SSIM \leq 1
\]

Higher SSIM means better perceptual quality.

\section{DCT Coefficient Visualization}

\begin{lstlisting}
def visualize_dct_coefficients(image):
    img_float = np.float32(image) - 128

    blocks = img_float.reshape(image.shape[0] // 8,
             8,
             image.shape[1] // 8,
             8).transpose(0, 2, 1, 3)

    dct_blocks = dctn(blocks, axes=(-2, -1), norm='ortho')

    avg_coeff = np.mean(np.abs(dct_blocks), axis=(0, 1))

    plt.figure(figsize=(5, 4))
    plt.imshow(np.log(avg_coeff + 1), cmap='hot')
\end{lstlisting}

\subsection{Purpose}

Visualizes average DCT coefficient magnitude across all image blocks.

\subsection{Why Use Absolute Value?}

DCT coefficients may be positive or negative.

Magnitude represents actual energy strength.

\subsection{Why Average Across Blocks?}

Shows global frequency distribution across the image.

\subsection{Why Log Scale?}

DCT coefficients vary significantly in magnitude.

Log scaling improves visibility of smaller coefficients.

\subsection{Interpretation}

Bright top-left region indicates:
\begin{itemize}
    \item strong low-frequency energy,
    \item dominant image structure.
\end{itemize}

Dark bottom-right region indicates:
\begin{itemize}
    \item weak high-frequency energy,
    \item less important fine details.
\end{itemize}

This validates the theoretical basis of JPEG compression.

\section{Image Loading and Preprocessing}

\begin{lstlisting}
image = cv2.imread("images/sample.jpg",
                   cv2.IMREAD_GRAYSCALE)
\end{lstlisting}

Loads grayscale image.

\subsection{Why Grayscale?}

Simplifies implementation and analysis.

\begin{lstlisting}
image = image[:h - h % 8, :w - w % 8]
\end{lstlisting}

Ensures dimensions are divisible by 8 for block processing.

\section{Compression Experiments}

\begin{lstlisting}
quality_levels = [5, 25, 50, 75, 95]
\end{lstlisting}

Multiple quality levels are tested to analyze the trade-off between:
\begin{itemize}
    \item compression ratio,
    \item reconstruction quality.
\end{itemize}

\section{Visualization of Results}

Compressed images are displayed at multiple quality levels.

\subsection{Expected Behavior}

\begin{itemize}
    \item Low quality $\rightarrow$ visible artifacts and high compression
    \item High quality $\rightarrow$ better reconstruction and larger size
\end{itemize}

\section{Graphs and Analysis}

PSNR and SSIM graphs visualize how image quality changes with compression level.

\subsection{Observations}

\begin{itemize}
    \item PSNR increases as quality increases.
    \item SSIM approaches 1 at high quality.
    \item Low quality factors introduce blocking artifacts.
\end{itemize}

\section{Conclusion}

This project successfully demonstrates the core principles of JPEG image compression using the Discrete Cosine Transform.

The experiment proves that:
\begin{itemize}
    \item DCT concentrates image energy into low frequencies.
    \item Quantization enables efficient lossy compression.
    \item High-frequency components can be removed with limited perceptual impact.
    \item PSNR and SSIM together provide meaningful image quality evaluation.
\end{itemize}

The implementation also demonstrates the efficiency of vectorized NumPy-based block processing compared to traditional nested loop implementations.

\end{document}